\chapter{Parent-origin энергии, состояния и скорости charged mediator}
\label{ch:v8-charged-mediator-energy-state-rate-origin}

\section{Симметрии оставляют одну свободную щель}

На минимальном Real-замкнутом environment-носителе
$\{|0\rangle,|+\rangle,|-\rangle\}$ имеем
\begin{equation}
 Y_E=\operatorname{diag}(0,1,-1),
 \qquad J_E|+\rangle=|-\rangle.
 \label{eq:v8-mediator-origin-charge-real-data}
\end{equation}
Вещественный самосопряжённый Hamiltonian, коммутирующий с $Y_E$ и $J_E$,
обязан иметь форму
\begin{equation}
 H_E=\operatorname{diag}(\varepsilon_0,\varepsilon_c,\varepsilon_c).
 \label{eq:v8-mediator-origin-hamiltonian-class}
\end{equation}
Линейные ограничения имеют ранг четыре на шестимерном
$\operatorname{Sym}_3(\mathbb R)$, поэтому допустимый класс двумерен. После
удаления общего energy shift остаётся ровно одна щель:
\begin{equation}
 H_E\sim\Delta_E(P_++P_-),
 \qquad \Delta_E=\varepsilon_c-\varepsilon_0.
 \label{eq:v8-mediator-origin-essential-gap}
\end{equation}
Нейтральная линия является основным состоянием только при
\begin{equation}
 \Delta_E>0,
 \qquad\text{а резонанс требует}\qquad
 \Delta_E=7\gamma.
 \label{eq:v8-mediator-origin-ground-and-resonance}
\end{equation}
Ни gauge-, ни Real-, ни family-симметрия не выбирает величину или знак
$\Delta_E$.

\section{Инвариантное состояние образует отрезок}

Общий gauge- и Real-инвариантный density operator равен
\begin{equation}
 \rho_E=\operatorname{diag}(p_0,p_c,p_c),
 \qquad p_0+2p_c=1,\qquad p_0,p_c\geq0.
 \label{eq:v8-mediator-origin-state-simplex}
\end{equation}
В частности, одинаково допустимы два разных состояния
\begin{equation}
 \rho_0=\operatorname{diag}(1,0,0),
 \qquad
 \rho_c=\operatorname{diag}(0,1/2,1/2).
 \label{eq:v8-mediator-origin-state-witnesses}
\end{equation}
Следовательно, симметрия не выбирает neutral vacuum. Если уже задана
положительная щель и $x=e^{-\beta\Delta_E}$, Gibbs-state имеет вид
\begin{equation}
 \rho_{\beta}=\frac{\operatorname{diag}(1,x,x)}{1+2x}.
 \label{eq:v8-mediator-origin-gibbs-state}
\end{equation}
При конечном $\beta\Delta_E$ он имеет ранг три, причём
\begin{equation}
 \Tr\rho_\beta^2=\frac{1+2x^2}{(1+2x)^2}<1,
 \qquad
 \det\rho_\beta=\frac{x^2}{(1+2x)^3}>0.
 \label{eq:v8-mediator-origin-gibbs-purity}
\end{equation}
Точный vacuum получается лишь как $\beta\Delta_E\to\infty$ или отдельной
чистой подготовкой.

\section{Нормировка connector не задаёт rate}

Следовая Gram-нормировка фиксирует форму двух квадратур, но interaction
\begin{equation}
 G_g=g\left(L_\downarrow\otimes b_+
 +L_\uparrow\otimes b_+^\dagger\right)
 \label{eq:v8-mediator-origin-free-coupling}
\end{equation}
допускает всякое $g>0$. Два свидетеля $g_1=1$ и $g_2=2$ сохраняют carrier,
симметрии и резонанс, но их weak-limit rates относятся как
\begin{equation}
 \frac{\Gamma_2}{\Gamma_1}=\frac{g_2^2}{g_1^2}=4.
 \label{eq:v8-mediator-origin-rate-witness}
\end{equation}
При физической clock-параметризации сохраняется лишь условная формула
\begin{equation}
 \Gamma=\frac{E_{\rm int}^2\tau_C}{\hbar^2}
 =\chi^2\frac{E_C}{\hbar}.
 \label{eq:v8-mediator-origin-clock-rate}
\end{equation}
Значения $E_C$ и $\chi$ ранее не были выведены.

\section{Существующая цепь не содержит новый канал}

Тёплиц-конвейер предоставляет переиспользуемый структурный шаблон
\begin{equation}
 V=(I\otimes S_{\rm chain})U_{\rm col}^{(0)},
 \qquad
 \Tr_{\rm chain}\operatorname{Ad}(V)^n
 (\rho\otimes\omega_{\rm vac})=\Phi_h^n(\rho).
 \label{eq:v8-mediator-origin-chain-template}
\end{equation}
Но его текущая ячейка содержит вакуум и $42$ старые jump-метки. Новый
singlet connector не входит в этот frame:
\begin{equation}
 \mathcal K_{\rm cell}^{\rm old}=\mathbb C|0\rangle\oplus\mathbb C^{42},
 \qquad b_+^{(s)}\notin\mathcal F_{42}.
 \label{eq:v8-mediator-origin-old-chain-mismatch}
\end{equation}
Поэтому требуется новая одинаковая cell-extension и заранее приготовленный
product vacuum. Их parent-origin не следует из существования шаблона.

Условная форма динамических данных закрыта:
\begin{equation}
 N_{\rm conditional\ shape}=8/8.
 \label{eq:v8-mediator-origin-conditional-ledger}
\end{equation}
Но пять физических входов остаются невыведенными:
\begin{equation}
 N_{\rm parent\ origin}=0/5:
 \quad\{\Delta_E=7\gamma,\rho_E=|0\rangle\langle0|,
 g,\text{new chain cell},(E_C,\chi)\}.
 \label{eq:v8-mediator-origin-parent-ledger}
\end{equation}

\begin{theorem}[Однопараметрический Hamiltonian без динамического origin]
Gauge- и Real-симметрии сокращают допустимый environment-Hamiltonian до
одной существенной щели, но не выбирают её знак или значение. Инвариантные
states образуют отрезок, connector допускает свободный coupling, а старый
42-channel conveyor не содержит новый jump. Поэтому energy/state/rate
parent-origin равен $0/5$.
\end{theorem}

\begin{nogocriterion}[Запрет сборки parent из условных слоёв]
Нельзя объявлять $\Delta_E=7\gamma$ следствием симметрии, neutral endpoint
--- вакуумом без state-selector, следовую нормировку --- физическим
coupling, а старый Toeplitz-шаблон --- уже подготовленной цепью нового
канала. Формула $\Gamma=\chi^2E_C/\hbar$ не фиксирует rate без обоих входов.
\end{nogocriterion}

\begin{protocol}\begin{enumerate}
\item допустимый Real/gauge Hamiltonian: \textbf{$\operatorname{diag}(\varepsilon_0,\varepsilon_c,\varepsilon_c)$};
\item существенных щелей: \textbf{$1$};
\item resonance выбран симметрией: \textbf{нет};
\item invariant state-space: \textbf{отрезок};
\item neutral vacuum выбран симметрией: \textbf{нет};
\item конечный Gibbs-state чист: \textbf{нет};
\item physical coupling выбран Gram-нормировкой: \textbf{нет};
\item rate witness $g_2/g_1=2$: \textbf{$\Gamma_2/\Gamma_1=4$};
\item старый 42-cell содержит новый connector: \textbf{нет};
\item conditional shape: \textbf{$8/8$};
\item energy/state/rate parent-origin: \textbf{$0/5$};
\item следующий гейт: \textbf{минимальная динамическая parent-архитектура}.
\end{enumerate}\end{protocol}

Машинный сертификат сохранён в
\path{s2t_v8_baryon_c0_singlet_triplet_central_gap_isotypic_channel_extra_edge_mass_two_source_cubic_incidence_up_down_binary_selector_minimal_typed_discriminator_binary_carrier_charged_environment_mediator_energy_state_rate_parent_origin_gate_results.json}.

\begin{thebibliography}{9}
\bibitem{v8-mediator-origin-attal} S.~Attal and Y.~Pautrat,
\emph{From Repeated to Continuous Quantum Interactions},
Ann. Henri Poincar\'e 7 (2006) 59; arXiv:math-ph/0311002.
\bibitem{v8-mediator-origin-bruneau} L.~Bruneau, A.~Joye and M.~Merkli,
\emph{Repeated Interactions in Open Quantum Systems},
J. Math. Phys. 55 (2014) 075204; arXiv:1305.2472.
\end{thebibliography}